\section{\texorpdfstring{$L$}{L}-functions}\label{sec:l-functions}
The characters of cyclic groups defined in \lemref{lem:chars-cyclic-group} are precisely the classical \emph{Dirichlet characters}.\index{Dirichlet, Peter}\index{L-function} Dirichlet characters are of paramount importance in analytic number theory: they appear in the study of primes in arithmetic progressions, the distribution of prime ideals, and the analytic continuation of $L$-functions. In the finite-field setting, the multiplicative characters of $\bbF_q$ play an analogous role, and the associated $L$-functions encode deep arithmetic information about character sums over $\bbF_q$.

Before specialising to the function-field setting, it is helpful to pause and describe the general analytic framework into which all of these objects fit. The central analytic device is a \emph{Dirichlet series}
\begin{Definition}{Dirichlet Series}{}
  \index{Dirichlet series}\index{L-function}%
  Let $a(n)$ be a sequence of complex numbers and $s\in\bbC$ be any complex number. Then a series of the form $$\sum_{n=1}^{\infty}\frac{a_n}{n^s}, \qquad s\in\bbC,$$ is called a Dirichlet series.
\end{Definition}
 Two very familiar Dirichlet series organize a vast section of analytic number theory:

\begin{enumerate}
  \item\textbf{Zeta functions} arise when the coefficients are the trivial ones, $a_n\equiv 1$, so that the Dirichlet series simply counts the size of some arithmetic object. The prototype is the \emph{Riemann zeta function} $\zeta(s)=\sum_n n^{-s}$, which encodes the multiplicative structure of $\bbZ$. 

  \item\textbf{$L$-functions} arise when the coefficients are \emph{twisted} by a character $\chi$, giving $L(s,\chi)=\sum_{h\colon\text{monic}} \chi(h)/\mfn(h)^s$. The character $\chi$ is what allows one to isolate arithmetic information that a zeta function averages over. Since $L$-functions are defined using polynomials over a finite fields and then we can talk about the prime analogue in polynomial rings i.e. irreducible polynomials. Taking the trivial character $\chi\equiv 1$ recovers $\zeta$ up to finitely many Euler factors, so zeta functions are the ``character-free'' special case of $L$-functions.
\end{enumerate}

The rest of this section unfolds this dictionary in two parallel worlds. First we recall the classical story over $\bbQ$ --- the Riemann zeta function, its zeros, and its $L$-function generalisations --- so that the relevant analytic phenomena are in view. Like in the case we go to $\bbQ$ the fraction field of $\bbZ$, we  develop the function-field $\bbF_q(X)$ analog: the zeta function associated to $\bbF_q[X]$ and its twists by multiplicative characters of $\bbF_q$. The zeta function turns out to be a rational function of $q^{-s}$, and the analog of the Riemann Hypothesis is a theorem (due to Weil) rather than an open problem.

\subsection{The Classical Riemann Zeta Function}
The prototype of everything that follows is the \emph{Riemann zeta function}.\index{Riemann, Bernhard}
\begin{Definition}{Riemann Zeta Function}{}
  \index{zeta function!Riemann}%
  For any complex number $s\in\bbC$, the zeta functions at $s$ is $$\zeta(s)=\sum_{n=1}^{\infty}\frac{1}{n^s}$$
\end{Definition}

which converges absolutely in the right half-plane $\Re(s)>1$. In fact for $s\in\bbC$ with $\Re(s)>1$ the product $\prod_{n=2}^{\infty}\zeta(n)$ also converges. Most proofs in this subsection are omitted, especially those relying on contour integration or the Mellin transform; the emphasis is on statements together with their arithmetic meaning.

\begin{Theorem}{Euler Productfor $\zeta(s)$}{euler-product-classical}%
  \index{Euler product!classical}\index{Euler product!function-field}%
  For $s\in\bbC$ with $\Re(s)>1$,
  $$\zeta(s)=\prod_{p\text{ prime}}\frac{1}{1-p^{-s}}.$$
\end{Theorem}
\begin{proof}
  Since $\bbZ$ is a UPD (unique factorization domain) every positive integer $n$ can be written uniquely as a product of powers of positive prime numbers. $$    \zeta(s) = \sum_{n=1}^{\infty}\frac{1}{n^s}= \prod_{p\colon\text{prime}}\lt(\sum_{i=0}^{\infty}\frac{1}{p^{i\cdot s}}\rt) = \prod_{p\colon\text{prime}} (1-p^{-s})^{-1}$$
  So we have the lemma.
\end{proof}
The identity is the analytic translation of \emph{unique factorisation}: expanding each geometric factor $(1-p^{-s})^{-1}=\sum_{k\geq 0}p^{-ks}$ and multiplying reproduces the sum over positive integers. Since $\zeta(s)$ has a pole at $s=1$ (the harmonic series diverges), the product must have infinitely many factors --- giving an analytic proof of the infinitude of primes. Dirichlet series evaluated on certain ``natural points'' encode interesting algebraic and arithmetic information.

The Euler product is the first of many important properties of $\zeta(s)$. The next is that $\zeta(s)$ has a meromorphic continuation\index{analytic continuation} to the entire plane; we recall the relevant definitions (zero, pole, meromorphic function) in Appendix~\ref{sec:meromorphic-analytic}. The special values of $\zeta$ encode deep arithmetic: $\zeta(2n)$ for $n \geq 1$ is a rational multiple of $\pi^{2n}$, expressible via Bernoulli numbers. Odd values are much harder --- Ap\'ery (1979) proved $\zeta(3)$ is irrational, while the arithmetic nature of $\zeta(5), \zeta(7), \ldots$ remains unknown. The negative-integer values $\zeta(0), \zeta(-1), \ldots$ make sense only \emph{after} analytic continuation: naively $\zeta(-1) = 1 + 2 + 3 + \cdots$, yet the continuation assigns it the value $-\tfrac{1}{12}$. To study $\zeta(s)$ and its analytic continuation we use the Gamma function $$\Gm(s)=\int_0^{\infty}e^{-t}t^{s-1}dt$$The analytic continuation is most cleanly stated for the \emph{completed zeta function}.

\begin{Theorem}{Analytic Continuation of the Completed Zeta Function}{zeta-functional-equation}%
  \index{zeta function!completed}\index{functional equation!for the zeta function}\index{functional equation}%
  Define the completed zeta function by
  $$\xi(s) = \tfrac{1}{2}s(s-1)\,\Gamma\!\lt(\tfrac{s}{2}\rt)\pi^{-s/2}\zeta(s).$$
  Then $\xi(s)$, originally defined for $\Re(s) > 1$, has an analytic continuation to an entire function and satisfies the \textbf{functional equation} $\xi(s) = \xi(1-s)$.
\end{Theorem}

From now on whenever we say zeta function we often mean the analytic continuation of zeta function by default. Now we will study the zeros of zeta function and why they are important. The Gamma function has poles at all negative integers. From the functional equation of the zeta function it has zeros at all the poles of $\Gm(\nicefrac{s}2)$ i.e. at all negative even integers. These zeros are called \emph{trivial zeros}\index{zeta function!trivial zeros} of $\zeta$. The $\zeta(s)$ has positive value for all $s\in\bbC$ with $\Re(s)>1$. Therefore all other zeros lie in the \emph{critical strip}\index{critical strip} $0\leq\Re(s)\leq 1$; the functional equation forces them to be symmetric under $s\mapsto 1-s$, and conjugation $s\mapsto\bar s$ forces reflection symmetry across the real axis. Hence 

\begin{Theorem}{Trivial Zeros and Symmetry of Non-trivial Zeros}{}
  The only zeros of $\zeta(s)$ where $s\in\bbC$ for $\Re(s)<0$ are at $s=-2n$ for $n\in\bbN$. In the critical strip the zeros lie symmetrically around the critical line $\Re(s)=\nicefrac12$ i.e. if $a+ib$ is a zero so is $1-a+ib$.
\end{Theorem}
regarding the zeros of the zeta functionThe line $\Re(s)=1/2$ is the \emph{critical line}\index{critical line}, and the most famous open problem in mathematics is:

\begin{Theoremcon}
  \index{Riemann hypothesis}\index{Riemann Hypothesis!classical}\index{Riemann Hypothesis!function fields}%
  \textcolor{mytheoremfr}{\textbf{Riemann Hypothesis}:} Every non-trivial zero of $\zeta(s)$ satisfies $\Re(s)=1/2$.
\end{Theoremcon}

While Hardy in 1914 showed that there are infinitely many zeros on the critical line the problem of every non-trivial zero being on the critical line is a major open problem in mathematics.
\begin{Theorem}{Hardy, 1914}{hardy}
  Infinitely many non-trivial zeros of $\zeta(s)$ lie on the critical line $\Re(s)=1/2$.
\end{Theorem}

Another interesting property, first noted in Riemann's original paper is the product expansion of zeta function using the non-trivial zeros.
\begin{Theorem}{Hadamard Product}{hadamard-product}%
  \index{Hadamard product}%
  There exist constants $A,B\in\bbC$ such that
  $$\xi(s)=e^{A+Bs}\prod_{\rho}\lt(1-\frac{s}{\rho}\rt)e^{s/\rho},$$
  where the product runs over all non-trivial zeros $\rho$ of $\zeta$.
\end{Theorem}

\begin{note}
  The \HadamardProdthm gives the ``zero-side'' factorisation of $\xi$, to be contrasted with the Euler product as the ``prime-side'' factorisation. Together they are the source of the famous duality between primes and zeros: equating the two factorisations --- typically after taking logarithmic derivatives --- converts statements about primes into statements about zeros and vice versa.
\end{note}

\paragraph{Zeros \texorpdfstring{$\leftrightarrow$}{<->} primes.} Taking the logarithmic derivative of the Euler product we obtain,
$$-\frac{\zeta'(s)}{\zeta(s)}=\sum_{n=1}^{\infty}\frac{\Lambda(n)}{n^s},$$
where $\Lambda(n)=\log p$ if $n=p^k$ and $0$ otherwise (the \emph{von Mangoldt function}\index{von Mangoldt function}). Therefore we can see how the non-trivial zeros of $\zeta(s)$ and the primes are very related.


\subsection{Dirichlet \texorpdfstring{$L$}{L}-function}%
\index{zeta function!over $\bbF_q[X]$}\index{zeta function!function-field}\index{prime polynomial}%
The function-field world has its own zeta function, obtained by replacing $\bbZ$ with the polynomial ring $\bbF_q[X]$ and the primes with monic irreducible polynomials. The notational parallel with the classical case is exact, but the analytic behavior is dramatically simpler: the resulting zeta function is a \emph{rational} function of $q^{-s}$, with no mystery in its analytic continuation, and (as will be developed in later sections) a fully-proven analog of the Riemann Hypothesis.

Let $\Phi\subseteq\bbF_q[X]$ denote the set of all monic polynomials, and let $\Phi_d\subseteq\Phi$ be the subset of polynomials of degree exactly $d$, for each $d\in\bbN$. A basic count gives $|\Phi_d|=q^d$: a degree-$d$ monic polynomial is determined by its $d$ lower-order coefficients, each of which is a free element of $\bbF_q$. This simple combinatorial fact is what makes the function-field zeta function so much more tractable than its classical counterpart.

For any polynomial $f\in\bbF_q[X]$, the quotient ring $\quotient{\bbF_q[X]}{\la f\ra}$ is a finite $\bbF_q$-algebra of size $q^{\deg(f)}$. We define the \emph{norm}\index{norm!of a polynomial} of $f$ to be
$$\mfn(f)=\lt|\quotient{\bbF_q[X]}{\la f\ra}\rt|=q^{\deg(f)},$$
by direct analogy with the norm of an ideal in the ring of integers of a number field. This norm is the direct replacement for the size of $\bbZ/n\bbZ$ (namely $n$) that appears in the classical Riemann zeta function.

To define $L$-functions in the function field setting we first need to extend multiplicative characters to the polynomial ring. Let $G\subseteq\bbF_q(X)$ be the group of rational functions $\nicefrac{h_1(X)}{h_2(X)}$ where $h_2\neq 0$ and $h_1,h_2\in\Phi$ where $\Phi\subseteq G$ is the set of all monic polynomials over $\bbF_q$. From now on we use $h(X)\in G$ to denote monic polynomials. Let $\ovG$ be the subgroup of $G$ such that $h_1\cdot h_2\in \ovG$ whenever $h_1,h_2\in\ovG$. For a character $\chi$ on $\ovG$ we extend $\chi$ to all of $G$ by setting $\chi(r)=0$ for all $r\notin \ovG$, after which $\chi$ remains multiplicative on $G$. This allows us to sum $\chi$ over all monic polynomials.
\begin{Definition}{Dirichlet $L$-function}{}
  \index{L-function!Dirichlet}\index{Dirichlet, Peter}%
  Let $\chi$ be a character on $\ovG$ extended to $G$ as above. For a complex number $s=a+ib\in\bbC$ with $a>1$, the \emph{Dirichlet $L$-function} associated to $\chi$ is defined by $$L(s,\chi)=\sum_{h\in\Phi}\chi(h)\cdot \mfn(h)^{-s}$$where the sum runs over all monic polynomials $h\in\Phi\subseteq G$ and $\mfn(h)=q^{\deg(h)}$ is the norm of $h$.
\end{Definition}

 Again this sum is absolutely convergent for $\Re(s)>1$. Like in the case of Riemann-Zeta functions we can express $L(s,\chi)$ in a Euler product form using only irreducible polynomials.
\begin{Theorem}{Euler Product for $\zeta(s)$}{thm:l-euler-product}%
  \index{Euler product!for $L$-functions over $\bbF_q$}%
  For $s\in\bbC$ with $\Re(s)>1$ we have $$L(s,\chi)=\prod_{\substack{h\in\Phi\\ h\text{ irreducible}}}(1-\chi(h)\cdot \mfn(h)^{-s})^{-1}$$
\end{Theorem}
\begin{proof}
  Since $\Re(s)>1$ we have \begin{align*}
    L(s,\chi) & = \sum_{h\in\Phi}\chi(h)\cdot \mfn(h)^{-s} = \prod_{\substack{h\in\Phi\\ h\text{ irreducible}}} \lt(\sum_{i=0}^{\infty} \chi(h^i)\cdot \mfn(h^i)^{-s} \rt) = \prod_{\substack{h\in\Phi\\ h\text{ irreducible}}}(1-\chi(h)\cdot \mfn(h)^{-s})^{-1}
  \end{align*}Hence we have the lemma.
\end{proof}
So let $\chi:G\to\bbC$ be a multiplicative complex valued function over the $G$ where $G$ is as defined above which also satisfies $|\chi(r)|\leq 1$ for all $r\in G$ and $\chi(1)=1$. Then for any $s\in\bbC$ with $\Re(s)>1$ we can rewrite $$L(s,\chi)=\sum_{h\in\Phi}\chi(h)\cdot \mfn(h)^{-s}=\sum_{d=0}^{\infty}\lt(\sum_{h\in\Phi_d}\chi(h)\rt)q^{-d\cdot s}$$
Since for any $h\in\Phi$ we have $\mfn(q)^s=\lt(q^s\rt)^{\deg(h)}$ we can consider the corresponding power series $$\ovL(z,\chi)=\sum_{h\in\Phi}\chi(h)\cdot z^{\deg(h)}$$hence $\ovL(q^{-s},\chi)=L(s,\chi)$. Now using \thmref{thm:l-euler-product} for $\ovL(z,\chi)$ we have $$\ovL(z,\chi)=\prod_{\substack{h\in\Phi\\ h\text{ irreducible}}}\lt(1-\chi(h)\cdot z^{\deg(h)}\rt)^{-1}$$Now we do some operations on the power series $\ovL(z,\chi)$ to get some expressions which help us later in proving some properties of exponential character sums. So for brevity of notations we use $\irr(\Phi)$ to denote the set of irreducible polynomials of $\Phi$. We have $\log \ovL(z,\chi)=-\sum\limits_{h\in\irr(\Phi)}\log \lt(1-\chi(h)z^{\deg(h)}\rt)$. So by differentiating we get $$\frac{d}{dz}\log \ovL(z,\chi) =-\sum\limits_{h\in\irr(\Phi)}\frac{-\chi(h)\cdot \deg(h)\cdot z^{\deg(h)-1}}{1-\chi(h)\cdot z^{\deg(h)}}=\sum\limits_{h\in\irr(\Phi)}\frac{\chi(h)\cdot \deg(h)\cdot z^{\deg(h)-1}}{1-\chi(h)\cdot z^{\deg(h)}}$$So \begin{align}
  z\cdot \frac{d}{dz}\log \ovL(z,\chi) & = \sum\limits_{h\in\irr(\Phi)}\frac{\chi(h)\cdot \deg(h)\cdot z^{\deg(h)-1}}{1-\chi(h)\cdot z^{\deg(h)}}\notag\\
  & = \sum_{h\in\irr(\Phi)}\chi(h)\cdot \deg(h)\cdot z^{\deg(h)}\sum_{k=0}^{\infty}\chi^k(h)\cdot z^{k\cdot \deg(h)}\notag\\
  & = \sum_{h\in\irr(\Phi)}\deg(h)\sum_{k=1}^{\infty}\chi^k(h)\cdot z^{k\cdot \deg(h)}\notag\\
  & = \sum_{k=1}^{\infty}\lt(\sum_{\substack{h\in\irr(\Phi),\\[0.5mm] \deg(h)\mid k}}\deg(h)\cdot \chi(h)^{\nicefrac{k}{\deg(h)}}\rt)z^k \label{eq:l-log-deriv-irreducibles}
\end{align}
So define $L_k\coloneqq  \sum_{h\in\irr(\Phi),\ \deg(h)\mid k}\deg(h)\cdot \chi(h)^{\nicefrac{k}{\deg(h)}}$. Hence we have $z\cdot \frac{d}{dz}\log \ovL(z,\chi)=\sum_{k=1}^{\infty}L_k\cdot z^k$. 

Now suppose there exists a positive integer $N\in\bbN$ such that for all $d>N$ we have $\sum_{h\in\Phi_d}\chi(g)=0$. Then we have $$L(s,\chi)=\sum_{d=0}^N \lt(\sum_{h\in\Phi_d}\chi(d)\rt)q^{-d\cdot s},\qquad  \ovL(z,\chi)=\sum_{d=0}^N \lt(\sum_{h\in\Phi_d}\chi(d)\rt)z^d$$
Since now $\ovL(z,\chi)$ is not an infinite power series anymore, instead a degree $N$ polynomial over $\bbC$ there exists complex numbers $w_1,\cdots, w_N\in\bbC$\index{reciprocal roots} such that $\ovL(z,\chi)=(1-w_1 z)(1-w_2 z)\cdots (1-w_N z)$.\index{L-function!finite degree} Now applying $z\cdot \frac{d}{dz}\log \ovL(z,\chi)$ on this expression we get \begin{align}
  z\cdot \frac{d}{dz}\log \ovL(z,\chi) & = -\sum_{i=1}^N\frac{w_iz}{1-w_iz}\notag\\
  & = -\sum_{i=1}^N w_iz\lt(\sum_{j=0}^{\infty}w_i^jz^j\rt)\notag\\
  & = -\sum_{i=1}^N\lt(\sum_{j=1}^{\infty}w_i^j\rt)z^i \label{eq:l-log-deriv-roots}
\end{align}
Therefore by comparing coefficients of \eqref{eq:l-log-deriv-roots} with \eqref{eq:l-log-deriv-irreducibles} we get $L_k=-\sum_{i=1}^{N}w_i^k$. for all $k\in\bbN$. So we have the following theorem
\begin{Theorem}{}{thm:l-function-root-factorisation}%
  \index{L-function!root factorisation}\index{L-function!factorisation}%
  Let $\chi:G\to\bbC$ be a multiplicative function with $|\chi(r)|\leq 1$ for all $r\in G$ and $\chi(1)=1$, and suppose there exists a positive integer $N\in\bbN$ such that $\sum_{h\in\Phi_d}\chi(h)=0$ for every $d>N$. Then $\ovL(z,\chi)$ is a polynomial of degree at most $N$ over $\bbC$, and there exist complex numbers $w_1,\dots,w_N\in\bbC$ such that
  $$\ovL(z,\chi)=\prod_{i=1}^{N}(1-w_iz), \qquad\text{ and }\qquad L(s,\chi)=\prod_{i=1}^{N}\lt(1-w_iq^{-s}\rt).$$
  Moreover, for every $k\in\bbN$,
  $$L_k\;=\;\sum_{\substack{h\in\irr(\Phi)\\ \deg(h)\mid k}}\deg(h)\cdot \chi(h)^{\nicefrac{k}{\deg(h)}}\;=\;-\sum_{i=1}^{N}w_i^k.$$
\end{Theorem}